\subsection{Программная реализация и особенности NEAT-модели}

Для проведения вычислительных экспериментов и моделирования когнитивных агентов в данной работе используется язык программирования Python и специализированная библиотека \cite{McIntyre_neat-python} (далее, \textit{neat-python}). Это удобный инструмент с открытым исходным кодом и строгим соответствием оригинальной архитектуре алгоритма NEAT, предложенной авторами работы \cite{stanley2002neuroevolution}.

Поведение алгоритма, параметры популяции и вероятности мутаций задаются централизованно через специальный конфигурационный файл. Ниже представлен анализ ключевых секций конфигурации, использовавшейся в ходе экспериментов, и их связь с теоретическим базисом нейроэволюции.

\subsubsection{Секция \texttt{[NEAT]}}

В процессе конфигурации алгоритма мы в первую очередь задали глобальные рамки. Размер популяции (pop\_size) был зафиксирован на протяжении всего эксперимента — это дало нам оптимальный баланс: генетического разнообразия хватало для поиска решений, а вычислительные мощности стенда не уходили в перегрузку. В качестве критерия останова был определен параметр fitness\_threshold. Как только лучший агент в симуляции пробивал этот потолок суммарной награды, задача считалась успешно решенной, и обучение автоматически прекращалось.

\subsubsection{Секция \texttt{[DefaultGenome]}}

Чтобы концепция инкрементального роста работала на практике, мы искусственно запретили алгоритму использовать скрытые слои на старте (num\_hidden = 0). Более того, параметр initial\_connection = full\_direct связывал сенсорные входы (num\_inputs) с моторными выходами (num\_outputs) напрямую на этапе инициализации агентов. В зависимости от задачи можно вырать необходимую функцию активации.

\subsubsection{Секции \texttt{[DefaultSpeciesSet]} и \texttt{[DefaultStagnation]}}

Далее нам потребовалось найти баланс между развитием новых топологий и изменения весов на заранее собранных. За темп роста графа отвечали вероятности структурных мутаций: добавление или удаление связей (add\_connection\_prob, delete\_connection\_prob) и самих узлов (add\_node\_prob, delete\_node\_prob). Параллельно с этим параметрические мутации (weight\_mutation\_prob, weight\_replace\_prob, weight\_perturb\_probs) обеспечивали изменение весов связей и узлов.

\subsubsection{Секции \texttt{[DefaultReproduction]} и \texttt{[DefaultStagnation]}}

Отдельного внимания потребовала настройка генетической дистанции $\delta$ (в соответствии с формулой \ref{eq:genetic_distance}). Чтобы защитить свежие структурные изменения, мы сделали явный упор на топологию: установили высокие штрафные коэффициенты за избыточные ($c_1$, параметр compatibility\_excess\_coefficient) и разъединенные гены ($c_2$, compatibility\_disjoint\_coefficient). Разница же в весовых коэффициентах ($c_3$, compatibility\_weight\_coefficient) наказывалась слабо, разрешая агентам внутри одного вида иметь разную силу синапсов. Границей, отделяющей один вид от другого, служил жесткий порог $\delta_t$ (compatibility\_threshold).Само формирование новых поколений шло по строгим правилам отбора. Мы задали процент выживших (survival\_threshold), пуская в размножение только наиболее результативных членов популяции. При этом параметр elitism гарантировал, что абсолютные чемпионы перейдут в следующую эпоху вообще без искажений от случайных мутаций. Чтобы эволюция не топталась на месте, механизмы species\_elitism и max\_stagnation отвечали за отсечение видов, надолго застрявших в развитии.